\subsubsection{Antecedentes de la investigación }

\cite{linocruz2023}, \textit{Implementación de una plataforma de chatbot gestionado por WhatsApp y Telegram para su comercialización en la empresa Innovus Software}, desarrollada en la Universidad Estatal del Sur de Manabí (UNESUM), en Jipijapa.

Este trabajo se centra en la mejora de la atención al cliente mediante un servicio de mensajería automática, orientado a responder consultas frecuentes de forma inmediata y precisa. A nivel metodológico, el estudio emplea revisión bibliográfica, análisis documental, enfoque histórico-lógico y observación del comportamiento de las consultas realizadas por los clientes.

Entre sus principales hallazgos, se reporta que una parte importante de las preguntas recibidas por la empresa son repetitivas y objetivas, lo que justifica la automatización del proceso de respuesta. La propuesta técnica se apoya en librerías (WhatsApp Cloud API) para la escucha de eventos de mensajería y en un algoritmo de secuencias programadas para identificar mensajes y generar respuestas automáticas.

Como resultado, la investigación concluye que la implementación de una plataforma de chatbot en canales como WhatsApp y Telegram contribuye a reducir costos operativos, incrementar la productividad y habilitar un servicio tecnológicamente viable para su comercialización a terceros. Este antecedente respalda la pertinencia del presente proyecto, especialmente en su enfoque de automatización de la comunicación para organizaciones, tiendas y negocios con alta demanda de atención digital.

\cite{vargasbraga_ifc}, en su trabajo \textit{Implementa{\c{c}}{\~a}o de assistente virtual via WhatsApp para aprimorar a comunica{\c{c}}{\~a}o no IFC -- Campus Sombrio}, presentan el desarrollo de un asistente virtual orientado a fortalecer la comunicaci{\'o}n institucional y la difusi{\'o}n de informaci{\'o}n oficial del campus. La propuesta se enfoca en responder dudas generales de los usuarios (por ejemplo, cursos, docentes y ubicaci{\'o}n), y se enmarca como una investigaci{\'o}n tecnol{\'o}gica basada en la construcci{\'o}n de un artefacto funcional. T{\'e}cnicamente, el sistema fue implementado con JavaScript/TypeScript e integrado con la API Baileys para la conexi{\'o}n con WhatsApp. Como resultado, reportan un asistente capaz de entregar respuestas simples de forma m{\'a}s r{\'a}pida que la b{\'u}squeda tradicional en redes sociales o en la web, destacando adem{\'a}s su potencial de ampliaci{\'o}n con nuevas funcionalidades.


\cite{zarabanda2024desarrollo} plantea el desarrollo de un chatbot interactivo para el proyecto curricular de Ingenier{\'i}a Telem{\'a}tica y Sistematizaci{\'o}n de Datos de la Universidad Distrital Francisco Jos{\'e} de Caldas, con el prop{\'o}sito de mejorar la comunicaci{\'o}n entre estudiantes y personal administrativo por medio de WhatsApp. La propuesta contempla funcionalidades orientadas a la gesti{\'o}n de tr{\'a}mites de radicaci{\'o}n de anteproyectos, notificaci{\'o}n de fechas clave, actualizaci{\'o}n de enlaces para documentaci{\'o}n requerida y atenci{\'o}n de preguntas frecuentes. A nivel arquitect{\'o}nico, el dise{\~n}o prioriza modularidad y escalabilidad: los mensajes son gestionados con Baileys y procesados en un servidor Node.js, que a su vez puede consultar MariaDB y APIs externas para construir respuestas adecuadas. Adem{\'a}s, el proyecto adopta metodolog{\'i}a SCRUM, con roles definidos (Product Owner, Scrum M{\'a}ster y equipo de desarrollo) y trabajo organizado en sprints con entregables progresivos desde la configuraci{\'o}n inicial hasta el despliegue y mantenimiento.

\cite{elsen2025hijab} desarrollan una aplicaci{\'o}n web de cat{\'a}logo electr{\'o}nico de productos hiyab integrada con un chatbot de WhatsApp para fortalecer la atenci{\'o}n al cliente en una MIPYME del sector moda musulmana en Indonesia. El estudio atiende problemas de gesti{\'o}n manual de inventario y baja capacidad de respuesta administrativa, proponiendo un sistema que permite consultar detalles de productos, precios, stock en tiempo real y promociones, adem{\'a}s de gestionar pedidos. El chatbot, implementado con la API Baileys, responde consultas frecuentes como ubicaci{\'o}n de tienda y valoraciones de productos, reduciendo carga operativa. La validaci{\'o}n reporta resultados satisfactorios de usabilidad (SUS 71.3 para no clientes y 73.5 para clientes), lo que respalda la pertinencia de integrar cat{\'a}logos web y agentes conversacionales para mejorar eficiencia operativa y experiencia de usuario.


\cite{vega2023nlp}, en su tesis de grado de la Universidad Mayor de San Andr{\'e}s, analiza el uso de procesamiento de lenguaje natural con modelos inteligentes para la atenci{\'o}n de consultas mediante redes sociales. El estudio describe la aplicaci{\'o}n de redes neuronales artificiales y arquitecturas Transformer para mejorar la comprensi{\'o}n del lenguaje en chatbots, con el objetivo de ofrecer interacciones m{\'a}s naturales y efectivas con los usuarios. Como caso aplicado, propone un chatbot para el gimnasio CrossFit capaz de informar sobre horarios, precios, actividades y recomendaciones de entrenamiento y nutrici{\'o}n. Entre los resultados comparativos, reporta mejor desempe{\~n}o del modelo de red neuronal artificial frente a GPT-2 y GPT-3 dentro de la metodolog{\'i}a empleada. Este antecedente aporta evidencia sobre la pertinencia de combinar t{\'e}cnicas de IA y PLN para fortalecer canales de atenci{\'o}n digital.
